%=============================================================================
% RELATED WORK SECTION
%=============================================================================

\paragraph{Spatial reasoning benchmarks.}
Evaluating spatial understanding in vision-language models has attracted considerable attention, yet existing benchmarks often conflate distinct cognitive abilities. The Visual Spatial Reasoning (VSR) dataset \cite{liu2023vsr} tests binary spatial relations between objects but relies on natural images where object recognition difficulty varies across samples. \citet{kamath2023whatsup} systematically probe VLMs on basic spatial relations and find consistent failures on vertical and horizontal judgments, though their evaluation focuses on absolute positions rather than relational reasoning. SpatialVLM \cite{chen2024spatialvlm} introduces metric distance estimation but does not decompose spatial reasoning into interpretable dimensions. Our benchmark addresses these limitations by isolating spatial reasoning from perceptual confounds through labeled visualizations, decomposing queries along egocentric-allocentric and viewer-object axes, and providing geometrically-computed ground truth that enables exact-match evaluation without annotation ambiguity.

\paragraph{Spatial cognition in vision-language models.}
Large vision-language models have demonstrated impressive capabilities across diverse tasks, yet their spatial reasoning abilities remain poorly characterized. GPT-4V \cite{openai2024gpt4} and Claude \cite{anthropic2024claude} exhibit emergent spatial understanding but lack systematic evaluation across spatial reasoning dimensions. Open-source models like Qwen2-VL \cite{wang2024qwen2vl} and LLaVA \cite{liu2024llava} have closed the gap on many vision-language tasks but show variable performance on spatial queries. Cognitive science distinguishes egocentric representations (viewer-centered) from allocentric representations (object-centered or world-centered) \cite{klatzky1998allocentric, burgess2006spatial}, with allocentric reasoning requiring mental transformation analogous to the classic mental rotation paradigm \cite{shepard1971mental}. Our benchmark operationalizes this distinction, enabling direct measurement of how well VLMs perform perspective transformation---a capability essential for embodied AI applications \cite{driess2023palme, brohan2023rt2} where agents must reason about spatial relations from viewpoints other than their own.

\paragraph{Mechanistic interpretability of visual reasoning.}
Understanding how neural networks implement reasoning has progressed substantially through mechanistic interpretability research. Work on transformer circuits \cite{elhage2021mathematical, olsson2022context} has revealed interpretable computational structures in language models, while recent scaling of these techniques to frontier models \cite{templeton2024scaling} demonstrates that meaningful features can be extracted from large-scale systems. For vision-language models, \citet{gandelsman2024interpreting} decompose CLIP representations via text-based analysis, revealing how visual concepts are encoded. However, mechanistic analysis of spatial reasoning in VLMs remains largely unexplored. Our work contributes to this direction by analyzing internal activations and attention patterns during spatial queries, identifying distinct failure mechanisms for egocentric versus allocentric tasks. This diagnostic approach moves beyond aggregate accuracy metrics to characterize how spatial reasoning succeeds or fails at the computational level.
